% Source: https://artofproblemsolving.com/community/c3h2016912_hard_problems
% Stable identity: aops-topic-2016912-q02
One digit is the average of the other two exactly when the three digits, in some order, form an arithmetic progression $p,\ p+d,\ p+2d$ with $d\ge0$.

\textbf{$d=0$.} All digits equal: $111,222,\ldots,999$, giving $9$ numbers.

\textbf{$d\ge1$.} The digits are distinct, with $0\le p$ and $p+2d\le9$. The number of such digit sets is $8,6,4,2$ for $d=1,2,3,4$, so $20$ sets. The $4$ sets with $p=0$ contain the digit $0$, which cannot lead, so each gives $3!-2!=4$ numbers; the other $16$ sets give $3!=6$ numbers each. This gives $4\cdot4+16\cdot6=112$ numbers.

So $9+112=121$ of the $900$ three-digit integers qualify, and the probability is $\boxed{\dfrac{121}{900}}$.
